%--------------------------------------------------------------------------
%	PACKAGES AND OTHER DOCUMENT CONFIGURATIONS
%--------------------------------------------------------------------------
\documentclass[12pt,a4paper]{article}
\usepackage[utf8]{inputenc}
\usepackage[francais]{babel}
\usepackage[T1]{fontenc}
\usepackage{amsmath}
\usepackage{amsfonts}
\usepackage{amssymb}
\usepackage{graphicx}
\usepackage{lmodern}
\usepackage[left=2cm,right=2cm,top=2.2cm,bottom=2cm]{geometry}

\usepackage{fancyhdr} % Required for custom headers
\usepackage{lastpage} % Required to determine the last page for the footer
\usepackage{extramarks} % Required for headers and footers
\usepackage[usenames,dvipsnames]{color} % Required for custom colors
\usepackage{graphicx} % Required to insert images
\usepackage{listings} % Required for insertion of code
\usepackage{courier} % Required for the courier font
\usepackage{verbatim}
\usepackage{multirow}

\usepackage{eurosym}

% Margins
%\topmargin=-0.45in
%\textwidth=6.5in
%\textheight=9.8in
\headsep=0.25in

% Set up the header and footer
%\pagestyle{fancy}
%\rhead{\firstxmark} % Top right header
%\lfoot{\lastxmark} % Bottom left footer
%\cfoot{} % Bottom center footer
%\rfoot{Page\ \thepage\ /\ \protect\pageref{LastPage}} % Bottom right footer
%\renewcommand\headrulewidth{0.3pt} % Size of the header rule
%\renewcommand\footrulewidth{0.3pt} % Size of the footer rule

%\setlength\parindent{0pt} % Removes all indentation from paragraphs

%--------------------------------------------------------------------------
%	CODE INCLUSION CONFIGURATION
%--------------------------------------------------------------------------

\definecolor{MyDarkGreen}{rgb}{0.0,0.4,0.0} % This is the color used for comments
\lstloadlanguages{C} % Load C syntax for listings, for a list of other languages supported see: ftp://ftp.tex.ac.uk/tex-archive/macros/latex/contrib/listings/listings.pdf

\begin{document}
	
%--------------------------------------------------------------------------
%	TITLE PAGE
%--------------------------------------------------------------------------
\begin{titlepage}
\newcommand{\HRule}{\rule{\linewidth}{0.5mm}} % Defines a new command for the horizontal lines, change thickness here
\centering % Center everything on the page
 
%	HEADING SECTIONS
\null
\vspace{2cm}
\textsc{\Large Université Catholique de Louvain}\\[1cm] % Name of your university/college
\textsc{\large LINGI1113 \\[0.3cm] Systèmes informatiques 2}\\[0.5cm] % Major heading such as course name
%\textsc{\large Minor Heading}\\[0.5cm] % Minor heading such as course title

%	TITLE SECTION

\HRule \\[0.4cm]
{ \LARGE \bfseries Projet 2~: PIC horloge\\[0.4cm] % Title of your document
\large \bfseries Rapport} \\[0.4cm]

\HRule \\[0.5cm]
 
\begin{figure}[!h]
	\begin{center}
	%2048 × 1364
		\includegraphics[width=10cm]{image.jpg}
	\end{center}
\end{figure}

%	AUTHOR SECTION

\large 
{\begin{tabular}{lll}
\textsc{Peschke} & Lena & 5826-11-00\\
\textsc{Sedda} & Mélanie & 2246-11-00\\
\end{tabular}}
\\[1cm]

\normalsize
{\begin{tabular}{ll}
\textit{Professeur} : & Marc Lobelle \\
\end{tabular}}
\\[0.7cm]

%	DATE SECTION

{\normalsize \today}\\[3cm] % Date, change the \today to a set date if you want to be precise

\newpage

\end{titlepage}

%--------------------------------------------------------------------------
%	TABLE OF CONTENTS
%--------------------------------------------------------------------------

\pagenumbering{gobble}
\clearpage
\thispagestyle{empty}
\tableofcontents
\clearpage
\pagenumbering{arabic}

%--------------------------------------------------------------------------
%	CONTENT
%--------------------------------------------------------------------------

\section{Documentation pour l'utilisateur}
Le réveil se compose de~:
\begin{itemize}
\item[1] écran LCD
\item[1] lampe LED jaune
\item[2] lampes LED rouges
\item[1] bouton MENU/NEXT/STOP (bouton 1)
\item[1] bouton SELECT/ADD/SNOOZE (bouton 2)
\end{itemize}

\subsection{Première utilisation}
Par défaut, l'horloge commence à 00~:~00~:~00, l'alarme est désactivée et réglée pour 00~:~00.
Lors de la mise sous tension du réveil, vous entrez dans le menu de l'horloge. Poussez SELECT si vous désirez changer l'heure, NEXT sinon.
Dans le premier cas, vous pouvez alors incrémenter les heures en appuyant sur ADD. Pour passer aux minutes et ensuite aux secondes, il suffit d'appuyer sur NEXT. Les crochets vous indiquent quelle valeur est en cours de modification.
Arrivé aux secondes, le bouton NEXT vous donne accès au menu de l'alarme, ou vous vous trouvez aussi si vous n'avez pas changé l'heure. Pour quitter le menu, poussez le bouton NEXT; pour régler l'alarme, poussez le bouton SELECT. En appuyant sur ADD, vous pouvez activer ou désactiver l'alarme. L'écran vous indique l'état actuel par ON ou OFF. Le bouton NEXT vous fait accéder à l'heure de l'alarme, dont vous pouvez changer la valeur en poussant sur ADD. NEXT vous permet de modifier les secondes avec le bouton ADD. Enfin, le bouton NEXT vous fait quitter le menu.

\subsection{L'affichage}
Le réveil affiche l'heure sous le format HH~:~MM~:~SS sur la première ligne de l'écran. La seconde indique si l'alarme est mise (Alarm ON) ou non (Alarm OFF) et à quelle heure elle est réglée. Pour accéder au menu, il suffit d'appuyer sur le bouton MENU et de le parcourir avec les boutons NEXT et SELECT. La LED jaune clignote en continu pour afficher le bon fonctionnement du réveil.

\subsection{Le réveil}
Lorsque le réveil sonne en faisant clignoter les deux LED rouges, il y a deux options possibles. Soit vous souhaitez l'éteindre; il suffit alors d'appuyer sur STOP. Le réveil sonnera à nouveau dans 24 heures. Soit vous désirez reporter le réveil de 5 minutes; appuyez alors sur SNOOZE. L'écran vous affiche maintenant l'heure en cours, le nombre de snobes effectués et l'heure de réveil d'origine. Pendant le mode snooze, ou à chaque fois que le réveil sonne à nouveau, vous pouvez reporter le réveil de 5 minutes supplémentaires, jusqu'à un maximum de 60 minutes au total. Vous pouvez aussi à tout moment appuyer sur STOP, ce qui aura pour effet de revenir à l'affichage normal de l'heure et de remettre l'alarme à l'heure d'origine pour une nouvelle utilisation le lendemain.

Il n'est pas possible d'accéder au menu tant que le réveil n'a pas été éteint. De même, tant que vous êtes dans le menu, le réveil ne sonne pas, pour ne pas interférer avec d'éventuels changements en cours.


\section{Documentation pour l'installateur}


\section{Documentation pour le programmeur}

\end{document}